% =====================================================================
%  BÁO CÁO BÀI TẬP LỚN - PHÂN TÍCH THIẾT KẾ HỆ THỐNG
%  Đề tài: XÂY DỰNG HỆ THỐNG QUẢN LÝ KHO HÀNG
%
%  Biên dịch (khuyến nghị):  pdflatex main.tex  (chạy 2 lần để cập nhật mục lục)
%  Yêu cầu gói: babel (vietnamese / vntex), tikz, listings, booktabs, ...
% =====================================================================
\documentclass[12pt,a4paper]{report}

% --- Hỗ trợ tiếng Việt: tương thích cả pdfLaTeX lẫn XeLaTeX/LuaLaTeX ---
\usepackage{iftex}
\ifPDFTeX
    \usepackage[utf8]{inputenc}
    \usepackage[T5]{fontenc}
    \usepackage[vietnamese]{babel}
\else
    \usepackage{fontspec}
    \usepackage{polyglossia}
    \setdefaultlanguage{vietnamese}
    \setmainfont{Latin Modern Roman}
\fi

\usepackage{lmodern}
\usepackage{geometry}
\geometry{a4paper, top=2.5cm, bottom=2.5cm, left=3cm, right=2cm}
\usepackage{graphicx}
\usepackage{booktabs}
\usepackage{array}
\usepackage{longtable}
\usepackage{tabularx}
\usepackage{multirow}
\usepackage{enumitem}
\usepackage{xcolor}
\usepackage{titlesec}
\usepackage{fancyhdr}
\setlength{\headheight}{14pt}
\usepackage{setspace}
\usepackage{caption}
\usepackage{float}
\usepackage{listings}
\usepackage{tikz}
\usetikzlibrary{shapes.geometric, arrows.meta, positioning, fit, backgrounds, calc}
\usepackage[hidelinks]{hyperref}

\onehalfspacing

% --- Màu sắc ---
\definecolor{primary}{HTML}{1F6FEB}
\definecolor{darkblue}{HTML}{10243E}
\definecolor{codegray}{HTML}{F4F6FB}
\definecolor{kwcolor}{HTML}{0B5394}
\definecolor{strcolor}{HTML}{1A7F37}
\definecolor{cmtcolor}{HTML}{6B7785}

% --- Định dạng tiêu đề chương ---
\titleformat{\chapter}[block]
  {\normalfont\Large\bfseries\color{darkblue}}
  {CHƯƠNG \thechapter.}{0.6em}{\Large}
\titlespacing*{\chapter}{0pt}{0pt}{16pt}

% --- Header/Footer ---
\pagestyle{fancy}
\fancyhf{}
\fancyhead[L]{\small\itshape Hệ thống Quản lý Kho hàng}
\fancyhead[R]{\small\thepage}
\renewcommand{\headrulewidth}{0.4pt}

% --- Cấu hình listings cho mã Python ---
\lstdefinestyle{py}{
  language=Python,
  backgroundcolor=\color{codegray},
  basicstyle=\ttfamily\footnotesize,
  keywordstyle=\color{kwcolor}\bfseries,
  stringstyle=\color{strcolor},
  commentstyle=\color{cmtcolor}\itshape,
  numbers=left, numberstyle=\tiny\color{cmtcolor}, numbersep=8pt,
  frame=single, rulecolor=\color{primary},
  breaklines=true, showstringspaces=false, tabsize=4,
  captionpos=b, literate={đ}{{\dj}}1 {á}{{\'a}}1 {à}{{\`a}}1
}
\lstset{style=py}

% Lệnh tiện ích cho ô bảng đặc tả use case
\newcolumntype{L}[1]{>{\raggedright\arraybackslash}p{#1}}

% Hình người (actor) vẽ bằng TikZ, dùng trong sơ đồ use case
\newcommand{\actorfig}{%
  \begin{tikzpicture}[scale=0.32, line width=0.7pt]
    \draw (0,2) circle (0.45);              % đầu
    \draw (0,1.55) -- (0,0.2);              % thân
    \draw (-0.8,1.1) -- (0.8,1.1);          % tay
    \draw (0,0.2) -- (-0.6,-0.7);           % chân trái
    \draw (0,0.2) -- (0.6,-0.7);            % chân phải
  \end{tikzpicture}%
}

\begin{document}

% =====================================================================
%  TRANG BÌA
% =====================================================================
\begin{titlepage}
\centering
{\large ĐẠI HỌC BÁCH KHOA HÀ NỘI\par}
{\large TRUNG TÂM ĐÀO TẠO LIÊN TỤC\par}
\vspace{0.3cm}
{\large ------o0o------\par}
\vspace{2.5cm}

{\Huge\bfseries BÀI TẬP LỚN\par}
\vspace{1.5cm}
{\large\itshape Đề tài:\par}
\vspace{0.4cm}
{\LARGE\bfseries\color{darkblue} XÂY DỰNG HỆ THỐNG\\[0.3cm] QUẢN LÝ KHO HÀNG\par}
\vspace{2.5cm}

\begin{flushleft}
\hspace{2.5cm}
\begin{tabular}{ll}
\textbf{GVHD:} & ThS. Phạm Thị Phương Giang \\[4pt]
\textbf{Lớp:} & B2CQ-CNTT01-K69 \\[4pt]
\textbf{Nhóm 9:} & 202490077: Bạch Minh Quang \\
\end{tabular}
\end{flushleft}

\vfill
{\large Hà Nội, tháng 06 năm 2026\par}
\end{titlepage}

% =====================================================================
%  MỤC LỤC
% =====================================================================
\tableofcontents
\listoffigures
\listoftables
\clearpage

% =====================================================================
\chapter{KHẢO SÁT HIỆN TRẠNG}
% =====================================================================

\section{Tình hình thực tế}

Hiện nay, hoạt động quản lý tại nhiều kho hàng của các doanh nghiệp vừa và nhỏ
vẫn chủ yếu dựa trên phương pháp bán thủ công: ghi chép sổ sách, theo dõi tồn
kho bằng các bảng tính rời rạc (Excel) hoặc giấy tờ. Khi quy mô kho mở rộng, số
lượng mặt hàng đa dạng và tần suất nhập -- xuất tăng cao, phương pháp truyền
thống bắt đầu bộc lộ sự quá tải, cồng kềnh và thiếu chính xác.

Bên cạnh đó, việc thiếu công cụ thống kê và báo cáo tự động khiến công tác quản
trị gặp nhiều hạn chế. Ban quản lý khó nắm bắt chính xác lượng tồn kho thực tế,
giá trị hàng tồn cũng như xu hướng luân chuyển hàng hóa để lập kế hoạch nhập
hàng hợp lý. Các quyết định mua bổ sung, phân bổ vốn lưu động thường dựa trên
cảm tính thay vì dữ liệu thực, làm giảm hiệu quả vận hành và dễ dẫn tới ứ đọng
vốn hoặc thiếu hàng.

Sự thiếu hụt một hệ thống phần mềm quản lý tập trung dẫn tới hàng loạt điểm
nghẽn (\textit{pain points}) trong quy trình nghiệp vụ:

\begin{itemize}[leftmargin=1.4cm]
  \item \textbf{Dữ liệu rời rạc, dễ sai sót:} Thông tin về mặt hàng, nhà cung
  cấp và phiếu nhập/xuất không được liên kết. Nhập liệu và đối soát thủ công
  thường xuyên dẫn tới sai lệch, trùng lặp hoặc thất lạc dữ liệu.

  \item \textbf{Tiêu tốn thời gian, nhân lực:} Nhân viên kho mất nhiều thời gian
  để tìm vị trí lưu trữ của hàng hóa, kiểm tra số lượng tồn, hay rà soát sổ sách
  để lập báo cáo định kỳ.

  \item \textbf{Thiếu kiểm soát, dễ thất thoát:} Không có cơ chế cảnh báo tự
  động, kho khó theo dõi các mặt hàng sắp hết hoặc tồn quá lâu. Việc đối soát
  giá trị tồn kho và truy vết nguồn gốc giao dịch diễn ra chậm trễ.

  \item \textbf{Khó phân quyền và truy vết trách nhiệm:} Khi nhiều người cùng
  thao tác trên một bảng tính chung, hệ thống không ghi nhận được ai đã thực
  hiện giao dịch nào, gây khó khăn cho việc kiểm tra và quy trách nhiệm.
\end{itemize}

\section{Các nhiệm vụ cần số hóa}

Từ những hạn chế trên, việc xây dựng một \textbf{Hệ thống Quản lý Kho hàng} là
yêu cầu cấp thiết. Hệ thống số hóa sẽ giúp chuẩn hóa và tự động hóa toàn bộ quy
trình nghiệp vụ: bảo đảm tính toàn vẹn dữ liệu, giảm áp lực cho nhân viên, tối
ưu hóa việc kiểm soát tồn kho và cung cấp số liệu chính xác phục vụ ra quyết
định.

Phạm vi đề tài tập trung số hóa và tự động hóa các quy trình quản trị cốt lõi
của kho: từ khâu tiếp nhận, đăng ký tài khoản người dùng cho tới quản lý chi
tiết hồ sơ cá nhân và phân quyền. Song song đó, hệ thống số hóa toàn bộ danh mục
hàng hóa và nhà cung cấp, giúp việc phân loại, cập nhật và tra cứu diễn ra khoa
học, chính xác.

Bên cạnh nền tảng quản trị dữ liệu, đề tài đi sâu vào bài toán luân chuyển hàng
hóa thông qua hai nghiệp vụ giao dịch trực tiếp là nhập kho và xuất kho, đồng
thời lưu trữ chặt chẽ lịch sử nhập/xuất để hỗ trợ tra soát. Cuối cùng, toàn bộ
dữ liệu vận hành được tổng hợp qua chức năng thống kê và báo cáo, cung cấp cho
ban quản lý cái nhìn toàn cảnh về tình hình kho và mức độ luân chuyển hàng hóa.

\section{Mô tả các nghiệp vụ}

\subsection{Nghiệp vụ 1: Tiếp nhận và đăng ký tài khoản người dùng}
\textbf{Mục đích:} Số hóa và quản lý danh sách người dùng (nhân viên kho, quản
lý kho). Bảo đảm mỗi cá nhân có một định danh duy nhất để truy cập và sử dụng
dịch vụ của hệ thống.

\noindent\textbf{Tác nhân:} Quản lý kho (Manager), Nhân viên kho (Staff).

\noindent\textbf{Mô tả quy trình:}
\begin{enumerate}[leftmargin=1.4cm]
  \item Quản lý kho cung cấp thông tin tài khoản mới (tên đăng nhập, họ tên,
  email, phòng ban, chức vụ, điện thoại) cho người dùng.
  \item Hệ thống kiểm tra tính hợp lệ và duy nhất của tên đăng nhập.
  \item Sau khi xác thực, hệ thống tạo tài khoản, băm mật khẩu và phân quyền
  tương ứng (Quản lý kho hoặc Nhân viên kho).
\end{enumerate}

\subsection{Nghiệp vụ 2: Quản lý hồ sơ cá nhân và phân quyền}
\textbf{Mục đích:} Lưu trữ, cập nhật thông tin cá nhân của người dùng; cho phép
đổi mật khẩu; phân quyền truy cập các chức năng theo vai trò.

\noindent\textbf{Tác nhân:} Người dùng, Quản lý kho.

\noindent\textbf{Mô tả quy trình:} Người dùng đăng nhập để xem và cập nhật thông
tin cá nhân (email, phòng ban, chức vụ, điện thoại) hoặc đổi mật khẩu. Hệ thống
xác thực mật khẩu cũ trước khi cho phép đổi mật khẩu mới. Vai trò người dùng
quyết định những chức năng được hiển thị: nhân viên kho thực hiện nghiệp vụ
nhập/xuất và tra cứu; quản lý kho có thêm quyền xóa dữ liệu và xem báo cáo tài
chính tổng hợp.

\subsection{Nghiệp vụ 3: Quản lý danh mục hàng hóa}
\textbf{Mục đích:} Số hóa danh mục hàng hóa trong kho, giúp tìm kiếm, phân loại
và kiểm kê chính xác, nhanh chóng.

\noindent\textbf{Tác nhân:} Nhân viên kho, Quản lý kho.

\noindent\textbf{Mô tả quy trình:} Khi có mặt hàng mới, nhân viên nhập thông tin
chi tiết: mã SKU, tên hàng, danh mục, đơn vị tính, ngưỡng tồn tối thiểu, đơn
giá, vị trí lưu trữ và nhà cung cấp. Hệ thống kiểm tra tính duy nhất của mã SKU.
Nhân viên có thể cập nhật thông tin, đổi trạng thái (Đang kinh doanh / Ngừng
kinh doanh). Hệ thống tự động cảnh báo những mặt hàng có tồn kho nhỏ hơn hoặc
bằng ngưỡng tối thiểu.

\subsection{Nghiệp vụ 4: Quản lý nhà cung cấp}
\textbf{Mục đích:} Lưu trữ và quản lý thông tin các đối tác cung cấp hàng hóa,
phục vụ việc liên hệ và truy xuất nguồn gốc hàng nhập.

\noindent\textbf{Tác nhân:} Nhân viên kho, Quản lý kho.

\noindent\textbf{Mô tả quy trình:} Người dùng thêm/sửa thông tin nhà cung cấp
gồm mã, tên, người liên hệ, điện thoại, email, địa chỉ và trạng thái hợp tác.
Mỗi mặt hàng có thể được gắn với một nhà cung cấp để thuận tiện cho việc đặt
hàng bổ sung.

\subsection{Nghiệp vụ 5: Xử lý nhập kho}
\textbf{Mục đích:} Ghi nhận và kiểm soát quá trình nhập hàng hóa vào kho từ nhà
cung cấp.

\noindent\textbf{Tác nhân:} Nhân viên kho.

\noindent\textbf{Mô tả quy trình:}
\begin{enumerate}[leftmargin=1.4cm]
  \item Nhân viên chọn mặt hàng cần nhập, nhập số lượng, đơn giá nhập và nhà
  cung cấp.
  \item Hệ thống kiểm tra hợp lệ (số lượng $>$ 0, đơn giá $\geq$ 0).
  \item Trong cùng một giao dịch CSDL, hệ thống \textit{cộng} số lượng vào tồn
  kho của mặt hàng và sinh phiếu nhập (mã phiếu dạng \texttt{PN-YYYYMMDD-\#\#\#\#}),
  ghi nhận thời gian, người thực hiện và thành tiền.
\end{enumerate}

\subsection{Nghiệp vụ 6: Xử lý xuất kho}
\textbf{Mục đích:} Ghi nhận quá trình xuất hàng hóa ra khỏi kho và kiểm soát
tồn kho.

\noindent\textbf{Tác nhân:} Nhân viên kho.

\noindent\textbf{Mô tả quy trình:}
\begin{enumerate}[leftmargin=1.4cm]
  \item Nhân viên chọn mặt hàng cần xuất, nhập số lượng, đơn giá xuất và bên
  nhận.
  \item Hệ thống kiểm tra điều kiện: tồn kho hiện tại phải lớn hơn hoặc bằng số
  lượng yêu cầu xuất. Nếu không đủ, hệ thống từ chối và thông báo lỗi.
  \item Nếu hợp lệ, trong cùng một giao dịch CSDL, hệ thống \textit{trừ} số
  lượng khỏi tồn kho và sinh phiếu xuất (mã phiếu dạng
  \texttt{PX-YYYYMMDD-\#\#\#\#}).
\end{enumerate}

\subsection{Nghiệp vụ 7: Quản lý lịch sử nhập/xuất}
\textbf{Mục đích:} Lưu vết (\textit{log}) toàn bộ giao dịch của từng mặt hàng và
từng người dùng để đối soát, truy xuất nguồn gốc và giải quyết khiếu nại.

\noindent\textbf{Tác nhân:} Nhân viên kho, Quản lý kho.

\noindent\textbf{Mô tả quy trình:} Sau mỗi nghiệp vụ nhập (NV5) hoặc xuất (NV6),
hệ thống tự động đẩy dữ liệu giao dịch vào bảng lịch sử. Người dùng có thể lọc
theo loại giao dịch (nhập/xuất) hoặc khoảng thời gian để tra cứu, kiểm tra ai là
người thực hiện giao dịch của một mặt hàng cụ thể.

\subsection{Nghiệp vụ 8: Thống kê và báo cáo}
\textbf{Mục đích:} Cung cấp số liệu tổng hợp trực quan giúp ban quản lý nắm bắt
tình hình kho, từ đó ra quyết định nhập bổ sung hoặc điều chỉnh.

\noindent\textbf{Tác nhân:} Quản lý kho.

\noindent\textbf{Mô tả quy trình:} Người quản lý chọn khoảng thời gian thống kê.
Hệ thống truy xuất dữ liệu và hiển thị các chỉ số: tổng số phiếu nhập/xuất, tổng
giá trị nhập/xuất, top hàng hóa nhập -- xuất nhiều nhất, cảnh báo tồn thấp và
tổng giá trị tồn kho. Hệ thống hỗ trợ kết xuất báo cáo giao dịch và báo cáo tồn
kho ra tệp định dạng CSV (mở được bằng Excel).

% =====================================================================
\chapter{MÔ HÌNH HÓA YÊU CẦU}
% =====================================================================

\section{Sơ đồ Use case tổng quát}

Sơ đồ use case tổng quát mô tả các tác nhân và những nhóm chức năng chính mà hệ
thống cung cấp. Hệ thống có hai tác nhân chính: \textbf{Nhân viên kho} (thực
hiện nghiệp vụ vận hành hằng ngày) và \textbf{Quản lý kho} (kế thừa toàn bộ
quyền của nhân viên và bổ sung quyền quản trị: xóa dữ liệu, xem báo cáo thống
kê tổng hợp).

\begin{figure}[H]
\centering
\begin{tikzpicture}[
    actor/.style={align=center, font=\small},
    usecase/.style={draw, ellipse, minimum width=3.3cm, minimum height=0.95cm,
                    align=center, font=\footnotesize, fill=primary!8},
    >={Stealth[]}
]
% Tác nhân (vẽ stick figure)
\node[actor] (staff) at (0,0) {\actorfig\\[2pt]Nhân viên\\kho};
\node[actor] (manager) at (12,0) {\actorfig\\[2pt]Quản lý\\kho};

% Khung hệ thống
\begin{scope}[on background layer]
\node[draw, rounded corners, fit={(2.4,-5.2) (9.6,4.6)},
      label=above:{\bfseries Hệ thống Quản lý Kho hàng}, fill=gray!3] {};
\end{scope}

% Use cases
\node[usecase] (login)    at (6, 4)    {Đăng nhập};
\node[usecase] (profile)  at (6, 2.7)  {Quản lý hồ sơ cá nhân};
\node[usecase] (product)  at (6, 1.4)  {Quản lý hàng hóa};
\node[usecase] (supplier) at (6, 0.1)  {Quản lý nhà cung cấp};
\node[usecase] (import)   at (6, -1.2) {Nhập kho};
\node[usecase] (export)   at (6, -2.5) {Xuất kho};
\node[usecase] (history)  at (6, -3.8) {Xem lịch sử giao dịch};
\node[usecase] (report)   at (6, -5)   {Báo cáo \& thống kê};

% Liên kết nhân viên kho
\foreach \uc in {login, profile, product, supplier, import, export, history}
    \draw (staff) -- (\uc);

% Liên kết quản lý kho (kế thừa nhân viên + báo cáo)
\foreach \uc in {report, history, export, import, supplier, product, profile, login}
    \draw (manager) -- (\uc);
\end{tikzpicture}
\caption{Sơ đồ Use case tổng quát của hệ thống}
\label{fig:usecase-tongquat}
\end{figure}

\noindent\textit{Ghi chú: hình người biểu diễn tác nhân (actor).
Quản lý kho có quan hệ kế thừa (generalization) với Nhân viên kho.}

\section{Sơ đồ Use case phân rã}

\subsection{Use case Quản lý hàng hóa}
Gói chức năng \textit{Quản lý hàng hóa} được phân rã thành các use case con:
\textit{Thêm hàng hóa}, \textit{Sửa hàng hóa}, \textit{Xóa hàng hóa} (chỉ Quản
lý kho), \textit{Tìm kiếm/Tra cứu} và \textit{Cảnh báo tồn thấp}. Use case
\textit{Thêm} và \textit{Sửa} đều \texttt{<<include>>} use case \textit{Kiểm tra
tính hợp lệ} (kiểm tra mã SKU duy nhất, dữ liệu bắt buộc).

\subsection{Use case Xử lý Nhập / Xuất kho}
Đây là nhóm use case thể hiện rõ quan hệ \texttt{<<include>>} và
\texttt{<<extend>>}. Use case \textit{Nhập kho} và \textit{Xuất kho} đều
\texttt{<<include>>} use case \textit{Ghi nhận giao dịch} và \textit{Cập nhật
tồn kho}. Use case \textit{Xuất kho} có điểm mở rộng \texttt{<<extend>>} là
\textit{Cảnh báo tồn không đủ} -- chỉ kích hoạt khi số lượng xuất vượt quá tồn
kho hiện có.

\begin{figure}[H]
\centering
\begin{tikzpicture}[
    actor/.style={align=center, font=\small},
    usecase/.style={draw, ellipse, minimum width=3cm, minimum height=0.9cm,
                    align=center, font=\footnotesize, fill=primary!8},
    >={Stealth[]}
]
\node[actor] (staff) at (0,-1) {\actorfig\\[2pt]Nhân viên\\kho};

\node[usecase] (import) at (4.5, 1)  {Nhập kho};
\node[usecase] (export) at (4.5, -3) {Xuất kho};
\node[usecase] (record) at (9, 0.2)  {Ghi nhận giao dịch};
\node[usecase] (update) at (9, -1.4) {Cập nhật tồn kho};
\node[usecase] (warn)   at (9, -3.4) {Cảnh báo tồn\\không đủ};

\draw (staff) -- (import);
\draw (staff) -- (export);

\draw[->, dashed] (import) -- node[above, font=\tiny, sloped] {<<include>>} (record);
\draw[->, dashed] (import) -- node[below, font=\tiny, sloped] {<<include>>} (update);
\draw[->, dashed] (export) -- node[above, font=\tiny, sloped] {<<include>>} (record);
\draw[->, dashed] (export) -- node[above, font=\tiny, sloped] {<<include>>} (update);
\draw[->, dashed] (warn)   -- node[above, font=\tiny, sloped] {<<extend>>} (export);
\end{tikzpicture}
\caption{Sơ đồ Use case phân rã: Xử lý Nhập / Xuất kho}
\label{fig:usecase-nhapxuat}
\end{figure}

\section{Đặc tả Use case}

\subsection{Use case 1: Đăng nhập hệ thống}
\begin{table}[H]
\centering
\renewcommand{\arraystretch}{1.3}
\begin{tabularx}{\textwidth}{|L{3.2cm}|X|}
\hline
\textbf{Tên Use case} & Đăng nhập hệ thống \\ \hline
\textbf{Tác nhân chính} & Nhân viên kho, Quản lý kho \\ \hline
\textbf{Mô tả tóm tắt} & Người dùng xác thực thông tin để truy cập vào các chức năng tương ứng với vai trò. \\ \hline
\textbf{Điều kiện} & Người dùng đã được cấp tài khoản. \\ \hline
\textbf{Luồng chính} &
1. Người dùng mở ứng dụng và nhập tên đăng nhập, mật khẩu. \newline
2. Hệ thống băm mật khẩu nhập vào và đối chiếu với CSDL (PBKDF2). \newline
3. Nếu sai: hiển thị thông báo lỗi. \newline
4. Nếu đúng: hệ thống xác định vai trò và điều hướng tới màn hình tương ứng. \\ \hline
\textbf{Ngoại lệ} & Tên đăng nhập hoặc mật khẩu không đúng $\rightarrow$ thông báo lỗi và yêu cầu nhập lại. \\ \hline
\end{tabularx}
\caption{Đặc tả use case Đăng nhập hệ thống}
\end{table}

\subsection{Use case 2: Đăng ký tài khoản người dùng}
\begin{table}[H]
\centering
\renewcommand{\arraystretch}{1.3}
\begin{tabularx}{\textwidth}{|L{3.2cm}|X|}
\hline
\textbf{Tên Use case} & Đăng ký tài khoản người dùng \\ \hline
\textbf{Tác nhân chính} & Quản lý kho \\ \hline
\textbf{Mô tả tóm tắt} & Quản lý kho tạo tài khoản mới cho nhân viên và phân quyền. \\ \hline
\textbf{Điều kiện} & Tài khoản chưa tồn tại trên hệ thống. \\ \hline
\textbf{Luồng chính} &
1. Quản lý nhập tên đăng nhập, mật khẩu, vai trò và thông tin hồ sơ. \newline
2. Hệ thống kiểm tra tính duy nhất của tên đăng nhập và độ dài mật khẩu. \newline
3. Nếu hợp lệ: hệ thống băm mật khẩu, tạo tài khoản và lưu vào CSDL. \\ \hline
\textbf{Ngoại lệ} & Tên đăng nhập đã tồn tại / mật khẩu quá ngắn $\rightarrow$ báo lỗi. \\ \hline
\end{tabularx}
\caption{Đặc tả use case Đăng ký tài khoản}
\end{table}

\subsection{Use case 3: Thêm hàng hóa mới}
\begin{table}[H]
\centering
\renewcommand{\arraystretch}{1.3}
\begin{tabularx}{\textwidth}{|L{3.2cm}|X|}
\hline
\textbf{Tên Use case} & Thêm hàng hóa mới \\ \hline
\textbf{Tác nhân chính} & Nhân viên kho, Quản lý kho \\ \hline
\textbf{Mô tả tóm tắt} & Người dùng nhập thông tin một mặt hàng mới vào danh mục kho. \\ \hline
\textbf{Điều kiện} & Người dùng đã đăng nhập. \\ \hline
\textbf{Luồng chính} &
1. Người dùng chọn ``Thêm hàng hóa'' và điền: SKU, tên, danh mục, đơn vị, tồn ban đầu, ngưỡng tối thiểu, đơn giá, vị trí, nhà cung cấp. \newline
2. Hệ thống kiểm tra hợp lệ (\texttt{<<include>>} Kiểm tra dữ liệu). \newline
3. Nếu mã SKU trùng: báo lỗi. \newline
4. Nếu hợp lệ: lưu mặt hàng vào CSDL và làm mới danh sách. \\ \hline
\textbf{Ngoại lệ} & Thiếu thông tin bắt buộc / SKU trùng $\rightarrow$ báo lỗi. \\ \hline
\end{tabularx}
\caption{Đặc tả use case Thêm hàng hóa}
\end{table}

\subsection{Use case 4: Xử lý nhập kho}
\begin{table}[H]
\centering
\renewcommand{\arraystretch}{1.3}
\begin{tabularx}{\textwidth}{|L{3.2cm}|X|}
\hline
\textbf{Tên Use case} & Xử lý nhập kho \\ \hline
\textbf{Tác nhân chính} & Nhân viên kho \\ \hline
\textbf{Mô tả tóm tắt} & Ghi nhận nhập hàng từ nhà cung cấp, tăng tồn kho và sinh phiếu nhập. \\ \hline
\textbf{Điều kiện} & Mặt hàng đã tồn tại trong danh mục. \\ \hline
\textbf{Luồng chính} &
1. Nhân viên chọn mặt hàng, nhập số lượng, đơn giá, nhà cung cấp. \newline
2. Hệ thống kiểm tra hợp lệ. \newline
3. Trong một giao dịch CSDL: cộng tồn kho \texttt{<<include>>} ghi phiếu nhập. \newline
4. Hệ thống thông báo mã phiếu và thành tiền. \\ \hline
\textbf{Ngoại lệ} & Số lượng/đơn giá không hợp lệ $\rightarrow$ hoàn tác (rollback) toàn bộ giao dịch. \\ \hline
\end{tabularx}
\caption{Đặc tả use case Xử lý nhập kho}
\end{table}

\subsection{Use case 5: Xử lý xuất kho}
\begin{table}[H]
\centering
\renewcommand{\arraystretch}{1.3}
\begin{tabularx}{\textwidth}{|L{3.2cm}|X|}
\hline
\textbf{Tên Use case} & Xử lý xuất kho \\ \hline
\textbf{Tác nhân chính} & Nhân viên kho \\ \hline
\textbf{Mô tả tóm tắt} & Ghi nhận xuất hàng ra khỏi kho, kiểm tra tồn và giảm tồn kho. \\ \hline
\textbf{Điều kiện} & Mặt hàng có tồn kho lớn hơn 0. \\ \hline
\textbf{Luồng chính} &
1. Nhân viên chọn mặt hàng, nhập số lượng, đơn giá, bên nhận. \newline
2. Hệ thống kiểm tra tồn kho $\geq$ số lượng xuất. \newline
3. Trong một giao dịch CSDL: trừ tồn kho \texttt{<<include>>} ghi phiếu xuất. \\ \hline
\textbf{Ngoại lệ (\texttt{<<extend>>})} & Tồn kho không đủ $\rightarrow$ hệ thống từ chối, hoàn tác giao dịch và cảnh báo. \\ \hline
\end{tabularx}
\caption{Đặc tả use case Xử lý xuất kho}
\end{table}

\subsection{Use case 6: Thống kê và báo cáo}
\begin{table}[H]
\centering
\renewcommand{\arraystretch}{1.3}
\begin{tabularx}{\textwidth}{|L{3.2cm}|X|}
\hline
\textbf{Tên Use case} & Thống kê và báo cáo \\ \hline
\textbf{Tác nhân chính} & Quản lý kho \\ \hline
\textbf{Mô tả tóm tắt} & Tổng hợp số liệu nhập/xuất, tồn kho trong khoảng thời gian và kết xuất báo cáo. \\ \hline
\textbf{Điều kiện} & Người dùng đăng nhập với vai trò Quản lý kho. \\ \hline
\textbf{Luồng chính} &
1. Quản lý chọn khoảng thời gian (Từ ngày -- Đến ngày). \newline
2. Hệ thống truy vấn và hiển thị: số phiếu nhập/xuất, tổng giá trị, top hàng hóa. \newline
3. Quản lý nhấn ``Xuất CSV''; hệ thống tạo tệp báo cáo và lưu xuống máy. \\ \hline
\textbf{Ngoại lệ} & Không có dữ liệu trong khoảng thời gian $\rightarrow$ hiển thị bảng rỗng. \\ \hline
\end{tabularx}
\caption{Đặc tả use case Thống kê và báo cáo}
\end{table}

\section{Biểu đồ lớp của ca sử dụng}

\subsection{Biểu đồ lớp cho ca sử dụng: Đăng nhập \& Đăng ký}
Biểu đồ thể hiện cấu trúc các lớp tham gia xác thực người dùng. Lớp giao diện
\texttt{WarehouseApp} gọi xuống \texttt{AccountService}; lớp này ủy quyền cho
\texttt{AuthService} xử lý trực tiếp với \texttt{UserRepository} để trả về hoặc
khởi tạo đối tượng \texttt{User}. Việc băm/so khớp mật khẩu do tiện ích
\texttt{security} (PBKDF2) đảm nhiệm.

\begin{figure}[H]
\centering
\begin{tikzpicture}[
  class/.style={draw, rectangle, fill=primary!6, align=left, font=\footnotesize,
                minimum width=2.6cm, inner sep=4pt},
  >={Stealth[]}, node distance=1.1cm and 1.6cm
]
\node[class] (app)  {\textbf{WarehouseApp}\\\rule{2.6cm}{0.3pt}\\+ login()\\+ register()};
\node[class] (acc)  [right=of app] {\textbf{AccountService}\\\rule{2.6cm}{0.3pt}\\+ login()\\+ register()\\+ updateProfile()};
\node[class] (auth) [right=of acc] {\textbf{AuthService}\\\rule{2.6cm}{0.3pt}\\+ login()\\+ register()\\+ changePassword()};
\node[class] (repo) [below=of auth] {\textbf{UserRepository}\\\rule{2.6cm}{0.3pt}\\+ findByUsername()\\+ insert()};
\node[class] (user) [left=of repo] {\textbf{User}\\\rule{2.6cm}{0.3pt}\\- username\\- role};
\draw[->] (app) -- (acc);
\draw[->] (acc) -- (auth);
\draw[->] (auth) -- (repo);
\draw[->] (repo) -- (user);
\end{tikzpicture}
\caption{Biểu đồ lớp ca sử dụng Đăng nhập \& Đăng ký}
\label{fig:class-auth}
\end{figure}

\subsection{Biểu đồ lớp cho ca sử dụng: Quản lý danh mục hàng hóa}
Luồng thêm/sửa/xóa/tìm kiếm hàng hóa. \texttt{CatalogService} chứa logic nghiệp
vụ (kiểm tra hợp lệ, kiểm tra trùng SKU), còn \texttt{ProductRepository} thực
hiện các câu lệnh SQL trên bảng \texttt{products} để truy xuất thực thể
\texttt{Product}.

\begin{figure}[H]
\centering
\begin{tikzpicture}[
  class/.style={draw, rectangle, fill=primary!6, align=left, font=\footnotesize,
                minimum width=2.6cm, inner sep=4pt},
  >={Stealth[]}, node distance=1.0cm and 1.7cm
]
\node[class] (app)  {\textbf{WarehouseApp}};
\node[class] (cat)  [right=of app] {\textbf{CatalogService}\\\rule{2.6cm}{0.3pt}\\+ addProduct()\\+ updateProduct()\\+ deleteProduct()\\+ searchProducts()};
\node[class] (prepo)[right=of cat] {\textbf{ProductRepository}\\\rule{2.6cm}{0.3pt}\\+ findAll()\\+ insert()\\+ update()\\+ deleteById()};
\node[class] (prod) [below=of prepo] {\textbf{Product}\\\rule{2.6cm}{0.3pt}\\- sku\\- quantity\\- unitPrice};
\node[class] (srepo)[below=of cat] {\textbf{SupplierRepository}};
\draw[->] (app) -- (cat);
\draw[->] (cat) -- (prepo);
\draw[->] (cat) -- (srepo);
\draw[->] (prepo) -- (prod);
\end{tikzpicture}
\caption{Biểu đồ lớp ca sử dụng Quản lý danh mục hàng hóa}
\label{fig:class-catalog}
\end{figure}

\subsection{Biểu đồ lớp cho ca sử dụng: Xử lý Nhập / Xuất kho}
Đây là ca sử dụng phức tạp nhất vì tác động đồng thời tới hai thực thể
(\texttt{Transaction} và \texttt{Product}). Logic nhập/xuất nằm ở
\texttt{InventoryService}; \texttt{TransactionRepository} sử dụng giao dịch CSDL
(\texttt{BEGIN}/\texttt{COMMIT}/\texttt{ROLLBACK}) để vừa ghi phiếu giao dịch,
vừa cập nhật số lượng tồn kho một cách an toàn.

\begin{figure}[H]
\centering
\begin{tikzpicture}[
  class/.style={draw, rectangle, fill=primary!6, align=left, font=\footnotesize,
                minimum width=2.6cm, inner sep=4pt},
  >={Stealth[]}, node distance=1.0cm and 1.7cm
]
\node[class] (app)  {\textbf{WarehouseApp}};
\node[class] (inv)  [right=of app] {\textbf{InventoryService}\\\rule{2.6cm}{0.3pt}\\+ importStock()\\+ exportStock()\\+ findAll()};
\node[class] (trepo)[right=of inv] {\textbf{TransactionRepository}\\\rule{2.6cm}{0.3pt}\\+ importStock()\\+ exportStock()\\\textit{(transaction CSDL)}};
\node[class] (tx)   [below=of trepo] {\textbf{Transaction}};
\node[class] (prod) [below=of inv] {\textbf{Product}};
\draw[->] (app) -- (inv);
\draw[->] (inv) -- (trepo);
\draw[->] (trepo) -- (tx);
\draw[->] (trepo) -- (prod);
\end{tikzpicture}
\caption{Biểu đồ lớp ca sử dụng Xử lý Nhập / Xuất kho}
\label{fig:class-inventory}
\end{figure}

\section{Biểu đồ trình tự}

\subsection{Nghiệp vụ Xuất kho}
Biểu đồ trình tự dưới đây mô tả luồng xử lý nghiệp vụ xuất kho -- nghiệp vụ điển
hình nhất vì có bước kiểm tra điều kiện và sử dụng giao dịch CSDL. Khi tồn kho
không đủ, hệ thống hoàn tác và trả về thông báo lỗi cho người dùng.

\begin{figure}[H]
\centering
\begin{tikzpicture}[font=\footnotesize, >={Stealth[]}]
  % Các lifeline
  \def\lifelines{6}
  \node (a1) at (0,0)  {\actorfig};
  \node (a2) at (3,0)  [draw, fill=primary!8] {:WarehouseApp};
  \node (a3) at (6.4,0)[draw, fill=primary!8] {:InventoryService};
  \node (a4) at (10,0) [draw, fill=primary!8] {:TransactionRepository};
  \node (a5) at (13.4,0)[draw, fill=primary!8] {:Database};
  \node[below=2pt of a1] {Nhân viên};

  % Đường lifeline dọc
  \foreach \n in {a1,a2,a3,a4,a5}
     \draw[dashed] (\n.south) -- ($(\n.south)+(0,-9)$);

  % Các thông điệp
  \draw[->] ($(a1.south)+(0,-0.8)$) -- node[above]{exportStock(...)} ($(a2.south)+(0,-0.8)$);
  \draw[->] ($(a2.south)+(0,-1.6)$) -- node[above]{exportStock(...)} ($(a3.south)+(0,-1.6)$);
  \draw[->] ($(a3.south)+(0,-2.4)$) -- node[above]{exportStock(...)} ($(a4.south)+(0,-2.4)$);
  \draw[->] ($(a4.south)+(0,-3.2)$) -- node[above]{BEGIN} ($(a5.south)+(0,-3.2)$);
  \draw[->] ($(a4.south)+(0,-4.0)$) -- node[above]{SELECT tồn kho} ($(a5.south)+(0,-4.0)$);
  \draw[<-] ($(a4.south)+(0,-4.8)$) -- node[above]{quantity} ($(a5.south)+(0,-4.8)$);

  % Khối alt
  \draw[thick] (5.6,-5.3) rectangle (14.2,-8.6);
  \node[anchor=west, fill=primary!15, font=\scriptsize] at (5.6,-5.3) {alt [tồn kho đủ / không đủ]};
  \draw[->] ($(a4.south)+(0,-5.9)$) -- node[above]{UPDATE trừ tồn} ($(a5.south)+(0,-5.9)$);
  \draw[->] ($(a4.south)+(0,-6.5)$) -- node[above]{INSERT phiếu xuất} ($(a5.south)+(0,-6.5)$);
  \draw[->] ($(a4.south)+(0,-7.1)$) -- node[above]{COMMIT} ($(a5.south)+(0,-7.1)$);
  \draw[dashed] (5.6,-7.5) -- (14.2,-7.5);
  \node[anchor=west, font=\scriptsize] at (5.7,-7.7) {[không đủ]};
  \draw[->] ($(a4.south)+(0,-8.2)$) -- node[above]{ROLLBACK} ($(a5.south)+(0,-8.2)$);

  % Trả kết quả
  \draw[<-, dashed] ($(a1.south)+(0,-8.9)$) -- node[above]{kết quả / thông báo} ($(a2.south)+(0,-8.9)$);
\end{tikzpicture}
\caption{Biểu đồ trình tự nghiệp vụ Xuất kho}
\label{fig:seq-export}
\end{figure}

\subsection{Các nghiệp vụ khác}
Các nghiệp vụ còn lại tuân theo cùng một khuôn mẫu trình tự phân tầng:
\textit{Tác nhân} $\rightarrow$ \texttt{WarehouseApp} (giao diện) $\rightarrow$
\textit{Service} (nghiệp vụ) $\rightarrow$ \textit{Repository} (truy cập dữ
liệu) $\rightarrow$ \texttt{Database}. Cụ thể:

\begin{itemize}[leftmargin=1.4cm]
  \item \textbf{Đăng nhập:} \texttt{WarehouseApp} $\to$ \texttt{AccountService}
  $\to$ \texttt{AuthService} $\to$ \texttt{UserRepository} (truy vấn người dùng,
  so khớp mật khẩu PBKDF2) $\to$ trả về \texttt{User}.
  \item \textbf{Nhập kho:} tương tự xuất kho nhưng \textit{cộng} tồn kho và bỏ
  qua bước kiểm tra điều kiện tồn.
  \item \textbf{Quản lý hàng hóa:} \texttt{WarehouseApp} $\to$
  \texttt{CatalogService} (kiểm tra hợp lệ) $\to$ \texttt{ProductRepository}
  $\to$ \texttt{Database}.
  \item \textbf{Thống kê -- báo cáo:} \texttt{WarehouseApp} $\to$
  \texttt{ReportService} (tổng hợp số liệu) $\to$ các Repository $\to$
  \texttt{Database}, sau đó kết xuất CSV.
\end{itemize}

% =====================================================================
\chapter{THIẾT KẾ HỆ THỐNG}
% =====================================================================

\section{Thiết kế lớp chi tiết}

Hệ thống được tổ chức theo \textbf{kiến trúc phân tầng} (layered architecture)
rõ ràng, mỗi tầng có trách nhiệm riêng và chỉ phụ thuộc vào tầng ngay bên dưới:

\begin{figure}[H]
\centering
\begin{tikzpicture}[
  layer/.style={draw, rounded corners, minimum width=12cm, minimum height=1.1cm,
                align=center, font=\small\bfseries}, >={Stealth[]}
]
\node[layer, fill=primary!20] (ui) at (0,0) {Tầng Giao diện (UI) -- \texttt{WarehouseApp} (Tkinter)};
\node[layer, fill=primary!14] (svc) at (0,-1.5) {Tầng Nghiệp vụ (Services)\\\footnotesize\mdseries AccountService, AuthService, CatalogService, InventoryService, ReportService};
\node[layer, fill=primary!9] (repo) at (0,-3.2) {Tầng Truy cập dữ liệu (Repositories)\\\footnotesize\mdseries UserRepository, ProductRepository, SupplierRepository, TransactionRepository};
\node[layer, fill=primary!6] (model) at (0,-4.9) {Tầng Mô hình miền (Models)\\\footnotesize\mdseries User, Product, Supplier, Transaction};
\node[layer, fill=gray!12] (infra) at (0,-6.4) {Hạ tầng \& Cấu hình -- \texttt{Database}, \texttt{DbConfig} (SQLite)};
\draw[->] (ui) -- (svc);
\draw[->] (svc) -- (repo);
\draw[->] (repo) -- (model);
\draw[->] (repo) -- (infra);
\end{tikzpicture}
\caption{Kiến trúc phân tầng tổng quan của hệ thống}
\label{fig:architecture}
\end{figure}

\subsection{Tầng giao diện}
Lớp \texttt{WarehouseApp} là lớp khởi động và điều phối UI trung tâm. Khi chạy,
phương thức \texttt{main()} khởi tạo \texttt{AppContext} (tạo \texttt{Database}
từ \texttt{DbConfig}, bảo đảm schema, nạp dữ liệu mẫu, khởi tạo các service),
sau đó dựng cửa sổ và hiển thị màn hình đăng nhập. Lớp này quản lý trạng thái
điều hướng (\texttt{current\_user}, màn hình hiện tại) và gọi các service tương
ứng khi người dùng thao tác.

\subsection{Tầng nghiệp vụ}
Gồm các lớp service xử lý logic chính:
\begin{itemize}[leftmargin=1.2cm]
  \item \texttt{AccountService} -- lớp trung gian cung cấp nghiệp vụ tài khoản,
  ủy quyền cho \texttt{AuthService}.
  \item \texttt{AuthService} -- xác thực đăng nhập, đăng ký, đổi mật khẩu (băm
  PBKDF2-HMAC-SHA256 có muối).
  \item \texttt{CatalogService} -- quản lý hàng hóa và nhà cung cấp, kiểm tra
  tính hợp lệ và trùng lặp dữ liệu.
  \item \texttt{InventoryService} -- điều phối nghiệp vụ nhập/xuất kho.
  \item \texttt{ReportService} -- tổng hợp số liệu thống kê và kết xuất báo cáo
  CSV.
\end{itemize}

\subsection{Tầng truy cập dữ liệu}
Gồm \texttt{UserRepository}, \texttt{ProductRepository},
\texttt{SupplierRepository} và \texttt{TransactionRepository}. Các repository
thực hiện câu lệnh SQL trực tiếp. Đặc biệt, \texttt{TransactionRepository} sử
dụng \textbf{giao dịch CSDL} cho hai thao tác \texttt{import\_stock} và
\texttt{export\_stock}: khi nhập kho, hệ thống kiểm tra hàng hóa, tăng số lượng
và thêm bản ghi giao dịch trong cùng một transaction; khi xuất kho, hệ thống
kiểm tra tồn kho, trừ số lượng và ghi giao dịch cũng trong cùng transaction,
bảo đảm tính toàn vẹn (atomicity).

Đoạn mã sau minh họa cơ chế giao dịch khi xuất kho:

\begin{lstlisting}[caption={Xử lý xuất kho với giao dịch CSDL (rút gọn)}]
def export_stock(self, *, product_id, quantity, unit_price, partner, username, note=""):
    with self._db.transaction() as conn:   # BEGIN ... COMMIT / ROLLBACK
        product = conn.execute(
            "SELECT id, name, quantity FROM products WHERE id = ?",
            (product_id,)).fetchone()
        if product is None:
            raise ProductNotFoundError(...)
        if product["quantity"] < quantity:          # kiem tra ton kho
            raise InsufficientStockError(...)        # -> rollback tu dong
        conn.execute("UPDATE products SET quantity = quantity - ? WHERE id = ?",
                     (quantity, product_id))         # tru ton kho
        code = self._next_code(conn, "EXPORT")       # sinh ma phieu
        conn.execute("INSERT INTO transactions (...) VALUES (...)", (...))
        return Transaction(...)
\end{lstlisting}

\section{Cơ sở dữ liệu}

Hệ thống sử dụng \textbf{SQLite} -- một hệ quản trị CSDL quan hệ nhúng, lưu toàn
bộ dữ liệu trong một tệp duy nhất, không cần cài đặt máy chủ. Lựa chọn này giúp
phần mềm chạy được ngay trên mọi máy có Python. Thiết kế lớp vẫn tách bạch phần
cấu hình (\texttt{DbConfig}) và kết nối (\texttt{Database}) nên có thể dễ dàng
chuyển sang MySQL/PostgreSQL khi triển khai thực tế.

\begin{table}[H]
\centering
\renewcommand{\arraystretch}{1.25}
\begin{tabularx}{\textwidth}{|L{4cm}|X|}
\hline
\textbf{Thuộc tính} & \textbf{Nội dung} \\ \hline
Hệ quản trị CSDL & SQLite 3 \\ \hline
Tên CSDL & \texttt{data/warehouse.db} \\ \hline
File schema chính & \texttt{sql/init.sql} (tự tạo trong \texttt{database.py}) \\ \hline
File cấu hình & \texttt{src/warehouse/config/db\_config.py} \\ \hline
Cơ chế bảo mật mật khẩu & PBKDF2-HMAC-SHA256 (có muối, 120.000 vòng lặp) \\ \hline
Mục đích & Quản lý hàng hóa, nhà cung cấp, người dùng và giao dịch nhập/xuất kho. \\ \hline
\end{tabularx}
\caption{Thông tin tổng quan cơ sở dữ liệu}
\end{table}

\noindent\textbf{Danh sách bảng:}
\begin{table}[H]
\centering
\renewcommand{\arraystretch}{1.2}
\begin{tabularx}{\textwidth}{|c|L{3.5cm}|X|}
\hline
\textbf{STT} & \textbf{Tên bảng} & \textbf{Mục đích} \\ \hline
1 & \texttt{users} & Tài khoản, mật khẩu mã hóa, vai trò và hồ sơ người dùng. \\ \hline
2 & \texttt{suppliers} & Thông tin nhà cung cấp hàng hóa. \\ \hline
3 & \texttt{products} & Danh mục hàng hóa và số lượng tồn kho. \\ \hline
4 & \texttt{transactions} & Lịch sử giao dịch nhập/xuất kho. \\ \hline
\end{tabularx}
\caption{Danh sách các bảng trong CSDL}
\end{table}

\subsection{Bảng \texttt{users}}
\begin{longtable}{|L{3cm}|L{3.2cm}|L{3.2cm}|L{4.4cm}|}
\hline
\textbf{Tên cột} & \textbf{Kiểu dữ liệu} & \textbf{Ràng buộc} & \textbf{Mô tả} \\ \hline
\endhead
id & INTEGER & PK, Auto Increment & Mã người dùng. \\ \hline
username & TEXT & Not Null, Unique & Tên đăng nhập. \\ \hline
password\_hash & TEXT & Not Null & Mật khẩu băm PBKDF2. \\ \hline
role & TEXT & Not Null, Check & Vai trò: \texttt{MANAGER} / \texttt{STAFF}. \\ \hline
full\_name & TEXT & Default '' & Họ và tên. \\ \hline
email & TEXT & Default '' & Email. \\ \hline
department & TEXT & Default '' & Phòng ban. \\ \hline
position & TEXT & Default '' & Chức vụ. \\ \hline
phone & TEXT & Default '' & Số điện thoại. \\ \hline
\caption{Cấu trúc bảng users}
\end{longtable}

\subsection{Bảng \texttt{suppliers}}
\begin{longtable}{|L{3cm}|L{3cm}|L{3.4cm}|L{4.4cm}|}
\hline
\textbf{Tên cột} & \textbf{Kiểu dữ liệu} & \textbf{Ràng buộc} & \textbf{Mô tả} \\ \hline
\endhead
id & INTEGER & PK, Auto Increment & Mã nhà cung cấp. \\ \hline
code & TEXT & Not Null, Unique & Mã định danh nhà cung cấp. \\ \hline
name & TEXT & Not Null & Tên nhà cung cấp. \\ \hline
contact\_person & TEXT & Default '' & Người liên hệ. \\ \hline
phone & TEXT & Default '' & Điện thoại. \\ \hline
email & TEXT & Default '' & Email. \\ \hline
address & TEXT & Default '' & Địa chỉ. \\ \hline
status & TEXT & Not Null, Check & \texttt{ACTIVE} / \texttt{INACTIVE}. \\ \hline
\caption{Cấu trúc bảng suppliers}
\end{longtable}

\subsection{Bảng \texttt{products}}
\begin{longtable}{|L{3cm}|L{2.8cm}|L{3.6cm}|L{4cm}|}
\hline
\textbf{Tên cột} & \textbf{Kiểu dữ liệu} & \textbf{Ràng buộc} & \textbf{Mô tả} \\ \hline
\endhead
id & INTEGER & PK, Auto Increment & Mã hàng hóa. \\ \hline
sku & TEXT & Not Null, Unique & Mã SKU. \\ \hline
name & TEXT & Not Null & Tên hàng hóa. \\ \hline
category & TEXT & Default '' & Danh mục. \\ \hline
unit & TEXT & Default 'cái' & Đơn vị tính. \\ \hline
quantity & INTEGER & Default 0, Check $\geq$ 0 & Số lượng tồn kho. \\ \hline
min\_quantity & INTEGER & Default 0, Check $\geq$ 0 & Ngưỡng tồn tối thiểu. \\ \hline
unit\_price & REAL & Default 0, Check $\geq$ 0 & Đơn giá. \\ \hline
location & TEXT & Default '' & Vị trí lưu trữ. \\ \hline
supplier\_id & INTEGER & Foreign Key & Tham chiếu \texttt{suppliers(id)}. \\ \hline
status & TEXT & Not Null, Check & \texttt{ACTIVE} / \texttt{DISCONTINUED}. \\ \hline
\caption{Cấu trúc bảng products}
\end{longtable}

\subsection{Bảng \texttt{transactions}}
\begin{longtable}{|L{3.2cm}|L{2.6cm}|L{3.2cm}|L{4cm}|}
\hline
\textbf{Tên cột} & \textbf{Kiểu dữ liệu} & \textbf{Ràng buộc} & \textbf{Mô tả} \\ \hline
\endhead
id & INTEGER & PK, Auto Increment & Mã giao dịch. \\ \hline
code & TEXT & Not Null, Unique & Mã phiếu (PN-/PX-...). \\ \hline
type & TEXT & Not Null, Check & \texttt{IMPORT} / \texttt{EXPORT}. \\ \hline
product\_id & INTEGER & Not Null, FK & Mã hàng hóa giao dịch. \\ \hline
product\_name & TEXT & Not Null & Tên hàng tại thời điểm giao dịch. \\ \hline
quantity & INTEGER & Check $>$ 0 & Số lượng. \\ \hline
unit\_price & REAL & Check $\geq$ 0 & Đơn giá giao dịch. \\ \hline
total & REAL & Check $\geq$ 0 & Thành tiền. \\ \hline
partner & TEXT & Default '' & Nhà cung cấp / bên nhận. \\ \hline
user\_username & TEXT & Not Null, FK & Người thực hiện. \\ \hline
transaction\_date & TIMESTAMP & Default now & Thời gian giao dịch. \\ \hline
note & TEXT & Default '' & Ghi chú. \\ \hline
\caption{Cấu trúc bảng transactions}
\end{longtable}

\section{Bảng sơ đồ quan hệ}

Mô hình thực thể -- liên kết (ERD) của hệ thống thể hiện bốn thực thể chính và
mối quan hệ giữa chúng:

\begin{figure}[H]
\centering
\begin{tikzpicture}[
  entity/.style={draw, rectangle, fill=primary!8, align=center, font=\small\bfseries,
                 minimum width=2.6cm, minimum height=1cm}, >={Stealth[]}
]
\node[entity] (sup)  at (0,0)    {suppliers};
\node[entity] (prod) at (5,0)    {products};
\node[entity] (tx)   at (10,0)   {transactions};
\node[entity] (user) at (10,-3)  {users};

\draw[->] (sup) -- node[above, font=\scriptsize]{1 \quad\quad N} (prod);
\draw[->] (prod) -- node[above, font=\scriptsize]{1 \quad\quad N} (tx);
\draw[->] (user) -- node[right, font=\scriptsize]{1 \;:\; N} (tx);
\end{tikzpicture}
\caption{Sơ đồ quan hệ thực thể (ERD) của hệ thống}
\label{fig:erd}
\end{figure}

\noindent\textbf{Quan hệ giữa các bảng:}
\begin{table}[H]
\centering
\renewcommand{\arraystretch}{1.2}
\begin{tabularx}{\textwidth}{|L{6cm}|X|}
\hline
\textbf{Quan hệ} & \textbf{Mô tả} \\ \hline
\texttt{suppliers.id} $\rightarrow$ \texttt{products.supplier\_id} & Một nhà cung cấp cung cấp nhiều mặt hàng. \\ \hline
\texttt{products.id} $\rightarrow$ \texttt{transactions.product\_id} & Một mặt hàng xuất hiện trong nhiều giao dịch. \\ \hline
\texttt{users.username} $\rightarrow$ \texttt{transactions.user\_username} & Một người dùng thực hiện nhiều giao dịch. \\ \hline
\end{tabularx}
\caption{Mô tả quan hệ giữa các bảng}
\end{table}

\noindent\textbf{Ràng buộc và chỉ mục quan trọng:}
\begin{table}[H]
\centering
\renewcommand{\arraystretch}{1.2}
\begin{tabularx}{\textwidth}{|L{4.2cm}|L{2.4cm}|X|}
\hline
\textbf{Tên} & \textbf{Loại} & \textbf{Mô tả} \\ \hline
PRIMARY KEY & Key & Khóa chính của cả bốn bảng. \\ \hline
\texttt{sku}, \texttt{username}, \texttt{code} & Unique & Không trùng lặp định danh nghiệp vụ. \\ \hline
\texttt{chk\_quantity} & Check & Số lượng tồn kho không âm. \\ \hline
\texttt{fk\_supplier} & Foreign Key & \texttt{products} tham chiếu \texttt{suppliers}. \\ \hline
\texttt{fk\_tx\_product} & Foreign Key & \texttt{transactions} tham chiếu \texttt{products}. \\ \hline
\texttt{fk\_tx\_user} & Foreign Key & \texttt{transactions} tham chiếu \texttt{users}. \\ \hline
\texttt{idx\_tx\_product/type/date} & Index & Tối ưu truy vấn lịch sử và thống kê. \\ \hline
\end{tabularx}
\caption{Các ràng buộc và chỉ mục}
\end{table}

\section*{Kết luận}
\addcontentsline{toc}{chapter}{Kết luận}
Đề tài đã hoàn thành việc phân tích, thiết kế và xây dựng \textbf{Hệ thống Quản
lý Kho hàng} bằng ngôn ngữ Python với giao diện Tkinter và CSDL SQLite. Hệ thống
đáp ứng đầy đủ tám nghiệp vụ cốt lõi: quản lý tài khoản, hồ sơ cá nhân, danh mục
hàng hóa, nhà cung cấp, nhập kho, xuất kho, lịch sử giao dịch và thống kê báo
cáo. Phần mềm được tổ chức theo kiến trúc phân tầng rõ ràng, sử dụng giao dịch
CSDL để bảo đảm tính toàn vẹn dữ liệu khi luân chuyển hàng hóa, và được kiểm thử
tự động (17 ca kiểm thử đơn vị) nhằm bảo đảm tính đúng đắn của các nghiệp vụ.

Hướng phát triển tiếp theo có thể bao gồm: triển khai trên CSDL máy chủ
(MySQL/PostgreSQL) cho môi trường nhiều người dùng, bổ sung mã vạch/QR cho hàng
hóa, tích hợp biểu đồ trực quan và xuất báo cáo PDF, cũng như xây dựng phiên bản
web để truy cập từ xa.

\end{document}
